\section{Conclusion}
\label{sec:conclusion}

Vision token pruning in VLAs is not a uniform operation: its outcome depends on the interaction between connector architecture and task complexity. Evaluating six VLA models (450M to 7B parameters) across four LIBERO task suites, we find that connector coupling is the primary determinant of pruning robustness. Token-independent connectors (linear projections, MLP projectors) degrade proportionally to the fraction removed, while token-entangling connectors (pixel-shuffle folding, BART encoder) exhibit task-amplified collapse. SmolVLA's pixel-shuffle connector produces 0.0\% success rate when just 5\% of tokens are removed; X-VLA retains its full 99.2\% baseline on Object at $k{=}0.9$ but loses 27.8 percentage points on Goal at the same ratio. The coefficient of variation of post-connector token L2 norms captures this spectrum as a single-pass diagnostic: models with CV$\,{=}\,$0.055 show $15{\times}$ greater method sensitivity than those with CV$\,{\geq}\,$0.15, and this metric correlates with both pruning tolerance and the retention threshold at which method choice becomes critical. For architectures where all pruning strategies fail, Vision Token Re-encoder recovers pixel-shuffle VLAs from 0.0\% to over 60\% success rate at 75\% compression through learned re-encoding.

These findings carry practical implications: computing CV before deployment can guide compression strategy selection without costly closed-loop evaluation. Two directions remain open. Real-robot validation is needed to confirm that the connector-complexity interaction observed in simulation transfers to physical manipulation. Adaptive per-task pruning schedules, where the retention ratio adjusts to estimated task complexity, could further close the gap between token-independent and token-entangling architectures.
